\section{Определение линейного функционала. Теорема о представлении линейного функционала. Определение биортогонального базиса. Теорема о биортогональном базисе. Закон преобразования координат линейного функционала при изменении базиса}

\begin{definition}
Пусть $V$ --- линейное пространство над полем $F$. Отображение $f: V \to F$ называется \textbf{линейным функционалом} (или линейной формой), если для любых $x, y \in V$ и любого $\lambda \in F$ выполняются условия:
\begin{enumerate}
    \item $f(x + y) = f(x) + f(y)$ (аддитивность);
    \item $f(\lambda x) = \lambda f(x)$ (однородность).
\end{enumerate}
Множество всех линейных функционалов на $V$ образует линейное пространство $V^*$, называемое \textbf{сопряженным} (или двойственным) пространством.
\end{definition}

\begin{theorem}[О представлении линейного функционала]
Пусть $E = \{e_1, e_2, \dots, e_n\}$ --- базис пространства $V$. Любой линейный функционал $f \in V^*$ полностью определяется своими значениями на базисных векторах $d_i = f(e_i)$ и для любого вектора $x = \xi_1 e_1 + \dots + \xi_n e_n$ вычисляется по формуле:
\[
f(x) = \sum_{i= is 1}^n d_i \xi_i.
\]
\end{theorem}

\begin{proof}
В силу линейности функционала и единственности разложения вектора $x$ по базису $E$, имеем:
\[
f(x) = f\left(\sum_{i=1}^n \xi_i e_i\right) = \sum_{i=1}^n \xi_i f(e_i) = \sum_{i=1}^n d_i \xi_i.
\]
Таким образом, значения $d_i$ однозначно задают функционал на всем пространстве.
\hfill$\blacksquare$
\end{proof}

\begin{definition}
Система функционалов $G = \{g_1, g_2, \dots, g_n\}$ в сопряженном пространстве $V^*$ называется \textbf{биортогональной} к базису $E = \{e_1, e_2, \dots, e_n\}$ пространства $V$, если для любых $i, j \in \{1, \dots, n\}$ выполняется условие:
\[
g_i(e_j) = \delta_{ij} = \begin{cases} 
1, & \text{если } i = j, \\ 
0, & \text{if } i \neq j,
\end{cases}
\]
где $\delta_{ij}$ --- символ Кронекера.
\end{definition}

\begin{theorem}[О существовании биортогонального базиса]
Для любого базиса $E = \{e_1, \dots, e_n\}$ пространства $V$ существует единственный биортогональный ему базис $G = \{g_1, \dots, g_n\}$ в сопряженном пространстве $V^*$. При этом $\dim V^* = n$.
\end{theorem}

\begin{proof}
\begin{enumerate}
    \item \textbf{Построение:} Для каждого $i = 1, \dots, n$ определим отображение $g_i: V \to F$ следующим образом: если $x = \xi_1 e_1 + \dots + \xi_n e_n$, то положим $g_i(x) = \xi_i$ (то есть $g_i$ сопоставляет вектору его $i$-ю координату в базисе $E$). В силу теоремы о линейности координат, каждое $g_i$ является линейным функционалом. По построению, $g_i(e_j) = \delta_{ij}$, значит, система $G$ биортогональна к $E$.
    \item \textbf{Линейная независимость:} Пусть $\sum_{i=1}^n \lambda_i g_i = 0$ (нулевой функционал). Применим эту линейную комбинацию к вектору $e_j$:
    \[
    \left(\sum_{i=1}^n \lambda_i g_i\right)(e_j) = \sum_{i=1}^n \lambda_i g_i(e_j) = \sum_{i=1}^n \lambda_i \delta_{ij} = \lambda_j = 0.
    \]
    Это верно для любого $j$, следовательно, все $\lambda_j = 0$, и система $G$ линейно независима.
    \item \textbf{Полнота (базисность):} Пусть $f$ --- произвольный функционал из $V^*$. Положим $d_i = f(e_i)$. Рассмотрим функционал $h = \sum_{i=1}^n d_i g_i$. Для любого базисного вектора $e_j$ имеем $h(e_j) = \sum_{i=1}^n d_i \delta_{ij} = d_j = f(e_j)$. Поскольку функционалы совпадают на базисе, они совпадают везде: $f = h$. Значит, $G$ порождает $V^*$.
\end{enumerate}
Таким образом, $G$ --- базис $V^*$, и $\dim V^* = n$.
\hfill$\blacksquare$
\end{proof}

\begin{theorem}[О преобразовании координат функционала]
Пусть $E$ и $E'$ --- два базиса пространства $V$, связанные матрицей перехода $T$ ($E' = E \cdot T$). Если функционал $f$ имеет в базисе $E$ строку координат $F = (d_1, \dots, d_n)$, а в базисе $E'$ --- строку координат $F' = (d'_1, \dots, d'_n)$, то они связаны соотношением:
\[
F' = F \cdot T.
\]
\end{theorem}

\begin{proof}
Пусть $x$ --- произвольный вектор со столбцом координат $X$ в базисе $E$ и $X'$ в базисе $E'$. Известно, что $X = T X'$. 
Значение функционала можно вычислить как скалярное произведение строки его координат на столбец координат вектора:
\[
f(x) = F X \quad \text{и} \quad f(x) = F' X'.
\]
Подставим $X = T X'$ в первое равенство:
\[
f(x) = F (T X') = (F T) X'.
\]
Поскольку это равенство должно выполняться для любого столбца $X'$, строки коэффициентов должны совпадать. Следовательно, $F' = F T$.
\hfill$\blacksquare$
\end{proof}
